% German Version - Methodological Article
\documentclass[12pt,a4paper]{article}
\usepackage[utf8]{inputenc}
\usepackage[T1]{fontenc}
\usepackage{lmodern}
\usepackage{amsmath,amssymb}
\usepackage{graphicx}
\usepackage{xcolor}
\usepackage{hyperref}
\usepackage{geometry}
\geometry{a4paper, left=3cm, right=3cm, top=3cm, bottom=3cm}
\usepackage{setspace}
\onehalfspacing
\usepackage{parskip}
\usepackage[ngerman]{babel}
\usepackage{csquotes}
\usepackage{microtype}
\usepackage{booktabs}
\usepackage{longtable}
\usepackage{array}
\usepackage{float}
\usepackage{url}
\usepackage{natbib}

\title{\Huge\textbf{Die empirische Grammatik von Marktgesprächen} \\[2mm]
       \LARGE Eine neuro-symbolische Rekonstruktion}

\author{
  \large
  \begin{tabular}{c}
    Paul Koop
  \end{tabular}
}

\date{\large 1994--2026}

\begin{document}

\maketitle

\begin{abstract}
Dieser Beitrag rekonstruiert die empirische Grammatik von acht Verkaufsgesprächen,
die 1994 auf dem Aachener Marktplatz aufgenommen wurden, mittels der Algorithmisch
Rekursiven Sequenzanalyse (ARS). Die Analyse führt von den Roh-Transkripten über
Terminalzeichenketten, Übergangszählung und Wahrscheinlichkeitsinduktion zu einer
probabilistischen kontextfreien Grammatik (PCFG) mit empirisch optimierten
Übergangswahrscheinlichkeiten. Ich biete eine interpretative Analyse der gelernten
Wahrscheinlichkeiten und unterscheide zwischen \textit{konstitutiven Regeln}
(structurellen Beschränkungen, die Wohlgeformtheit definieren) und \textit{statistischen
Regularitäten} (empirischen Häufigkeiten, die kontingente Muster widerspiegeln).
Die Grammatik offenbart sowohl ritualisierte Sequenzen (Begrüßungen, Verabschiedungen)
als auch strategische Optionen (Upselling-Schleifen, Neustarts). Der Beitrag schließt
mit einer Diskussion, wie die Grammatik in neuro-symbolischen Frameworks implementiert
werden kann, um qualitative Interpretation und computergestützte Ausführung zu
verbinden.
\end{abstract}

\newpage
\tableofcontents
\newpage

\section{Einleitung: Von Rohdaten zur formalen Grammatik}

Die Algorithmisch Rekursive Sequenzanalyse (ARS) beruht auf einer einfachen,
aber wirkmächtigen Idee: Soziale Interaktionen hinterlassen physikalische Spuren
(Tonaufnahmen), die transkribiert, interpretiert und in formale Grammatiken
überführt werden können. Dieser Beitrag dokumentiert die gesamte Pipeline von
den Roh-Transkripten bis zur empirisch optimierten Grammatik, basierend auf
acht Verkaufsgesprächen, die im Juni/Juli 1994 auf dem Aachener Marktplatz
aufgenommen wurden.

Der Beitrag ist dreifach:

\begin{enumerate}
    \item \textbf{Empirisch}: Ich präsentiere die vollständige optimierte
    Grammatik, induziert aus den acht Transkripten, mit allen Übergangs-
    wahrscheinlichkeiten.
    
    \item \textbf{Interpretativ}: Ich analysiere die gelernten Wahrscheinlichkeiten
    und unterscheide zwischen konstitutiven Regeln (strukturellen Notwendigkeiten)
    und statistischen Regularitäten (empirischen Kontingenzen).
    
    \item \textbf{Methodologisch}: Ich diskutiere, wie die Grammatik in
    neuro-symbolischen Frameworks implementiert werden kann, um qualitative
    Interpretation und computergestützte Ausführung zu verbinden.
\end{enumerate}

\section{Daten und Interpretation}

\subsection{Die acht Transkripte}

Das empirische Material umfasst acht Transkripte von Verkaufsgesprächen, die
im Juni/Juli 1994 auf dem Aachener Marktplatz aufgenommen wurden. Die Transkripte
variieren in Länge, Vollständigkeit und Gesprächstyp:

\begin{table}[H]
\centering
\caption{Korpus der Marktgespräche}
\label{tab:corpus}
\begin{tabular}{@{} l l l l @{}}
\toprule
\textbf{Text} & \textbf{Datum} & \textbf{Ort} & \textbf{Beteiligte} \\
\midrule
T1 & 28.06.1994 & Metzgerei & Verkäuferin (w), Kunde \\
T2 & 28.06.1994 & Kirschenstand & Verkäufer (m), K1, K2 \\
T3 & 28.06.1994 & Fischstand & Verkäufer (m), Kunde \\
T4 & 28.06.1994 & Gemüsestand & Verkäufer (m), Kunde \\
T5 & 26.06.1994 & Gemüsestand & Verkäufer (m), K1, K2 \\
T6 & 28.06.1994 & Käseverkauf & Verkäufer (m), K1 \\
T7 & 28.06.1994 & Bonbonstand & Verkäufer (m), Kunde \\
T8 & 09.07.1994 & Bäckerei & Verkäuferin (w), Kunde \\
\bottomrule
\end{tabular}
\end{table}

\subsection{Zuordnung der Terminalzeichen}

Jede Äußerung wurde einem Terminalzeichen aus einem vordefinierten Kategoriensystem
zugeordnet, das durch sequenzielle Mikroanalyse nach der Methode der objektiven
Hermeneutik entwickelt wurde \citep{oevermann1979methodology}. Die zwölf
Terminalzeichen sind:

\begin{table}[H]
\centering
\caption{Terminalzeichen und ihre Bedeutungen}
\label{tab:terminals}
\begin{tabular}{@{} l l @{}}
\toprule
\textbf{Symbol} & \textbf{Bedeutung} \\
\midrule
KBG & Kunden-Gruß \\
VBG & Verkäufer-Gruß \\
KBBd & Kunden-Bedarf (konkret) \\
VBBd & Verkäufer-Nachfrage \\
KBA & Kunden-Antwort \\
VBA & Verkäufer-Reaktion \\
KAE & Kunden-Erkundigung \\
VAE & Verkäufer-Auskunft \\
KAA & Kunden-Abschluss \\
VAA & Verkäufer-Abschluss \\
KAV & Kunden-Verabschiedung \\
VAV & Verkäufer-Verabschiedung \\
\bottomrule
\end{tabular}
\end{table}

Die resultierenden Terminalzeichenketten für die acht Transkripte sind:

\begin{verbatim}
T1: KBG, VBG, KBBd, VBBd, KBA, VBA, KBBd, VBBd, KBA, VAA, KAA, VAV, KAV
T2: VBG, KBBd, VBBd, VAA, KAA, VBG, KBBd, VAA, KAA
T3: KBBd, VBBd, VAA, KAA
T4: KBBd, VBBd, KBA, VBA, KBBd, VBA, KAE, VAE, KAA, VAV, KAV
T5: KAV, KBBd, VBBd, KBBd, VAA, KAV
T6: KBG, VBG, KBBd, VBBd, KAA
T7: KBBd, VBBd, KBA, VAA, KAA
T8: KBG, VBBd, KBBd, VBA, VAA, KAA, VAV, KAV
\end{verbatim}

\section{Die optimierte Grammatik}

\subsection{Induktionsmethode}

Die Grammatik wurde durch Zählung der Übergänge über alle acht Transkripte
induziert und zu Wahrscheinlichkeiten normalisiert:

\[
P(\sigma_j \mid \sigma_i) = \frac{\text{Zahl}(\sigma_i \to \sigma_j)}{\sum_k \text{Zahl}(\sigma_i \to \sigma_k)}
\]

Die resultierenden Wahrscheinlichkeiten wurden iterativ verfeinert, indem
künstliche Ketten generiert und deren Häufigkeitsverteilungen mit den empirischen
Daten verglichen wurden. Die finale Grammatik erreichte eine Korrelation von
\(r = 0,925\) zwischen empirischen und generierten Häufigkeiten.

\subsection{Die optimierte probabilistische Grammatik}

\begin{table}[H]
\centering
\caption{Optimierte Übergangswahrscheinlichkeiten}
\label{tab:grammar}
\begin{tabular}{@{} l l @{}}
\toprule
\textbf{Startsymbol} & \textbf{Folgesymbole mit Wahrscheinlichkeiten} \\
\midrule
KBG & VBG (0,667), VBBd (0,333) \\
VBG & KBBd (1,0) \\
KBBd & VBBd (0,667), VAA (0,167), VBA (0,167) \\
VBBd & KBA (0,444), VAA (0,222), KBBd (0,222), KAA (0,111) \\
KBA & VBA (0,5), VAA (0,5) \\
VBA & KBBd (0,5), KAE (0,25), VAA (0,25) \\
VAA & KAA (0,857), KAV (0,143) \\
KAA & VAV (0,75), VBG (0,25) \\
VAV & KAV (1,0) \\
KAE & VAE (1,0) \\
VAE & KAA (1,0) \\
KAV & VAV (0,5), KBBd (0,5) \\
\bottomrule
\end{tabular}
\end{table}

\subsection{Grafische Darstellung}

Abbildung~\ref{fig:grammar} visualisiert die Grammatik als gerichteten Graphen
mit Wahrscheinlichkeiten.

\begin{figure}[H]
\centering
\begin{verbatim}
                    KBG
                   /   \
                0,667   0,333
                 /       \
               VBG       VBBd
                |          |
               1,0         |
                |          |
               KBBd        |
               / | \        |
           0,667|  |0,167   |
             /  |  \        |
           VBBd |   VBA     |
            |   |    |      |
            |   |  0,5|     |
            |   |    |      |
            |   |   KBBd    |
            |   |   / \     |
            |   |  ... ...  |
            |   |           |
            \---/-----------/
               |
              VAA
             /   \
          0,857   0,143
           /       \
         KAA       KAV
         / \        |
      0,75 0,25     |
       /    \       |
     VAV    VBG     |
      |             |
     KAV           VAV
      |             |
      \_____________/
            |
           ENDE
\end{verbatim}
\caption{Grafische Darstellung der optimierten Grammatik. Kantengewichte
geben die Übergangswahrscheinlichkeiten an.}
\label{fig:grammar}
\end{figure}

\section{Interpretative Analyse der Wahrscheinlichkeiten}

\subsection{Konstitutive Regeln (strukturelle Notwendigkeiten)}

Einige Übergänge haben die Wahrscheinlichkeit 1,0, was darauf hinweist, dass
sie nicht bloß statistische Regularitäten, sondern \textit{konstitutive Regeln}
des Interaktionsformats sind:

\begin{itemize}
    \item \textbf{VBG → KBBd (1,0)}: Ein Verkäufer-Gruß muss von einer
    Kunden-Bedarfsäußerung gefolgt sein. Dies ist eine strukturelle Eigenschaft
    von Verkaufsgesprächen. Ohne diesen Übergang würde dem Gespräch sein
    transaktionaler Zweck fehlen.
    
    \item \textbf{KAE → VAE (1,0)}: Eine Kunden-Erkundigung muss von einer
    Verkäufer-Auskunft beantwortet werden. Dies spiegelt die normative
    Erwartung von Reziprozität und Kooperation wider.
    
    \item \textbf{VAE → KAA (1,0)}: Auf Verkäufer-Auskunft folgt immer ein
    Kunden-Abschluss. Dies deutet darauf hin, dass der Informationsaustausch
    in Marktgesprächen eng mit der Transaktionsbeendigung gekoppelt ist.
    
    \item \textbf{VAV → KAV (1,0)}: Verkäufer-Verabschiedungen werden immer
    durch Kunden-Verabschiedungen erwidert. Dies ist eine ritualisierte
    Abschlusssequenz, die das Ende der Interaktion markiert.
\end{itemize}

\subsection{Statistische Regularitäten (empirische Kontingenzen)}

Andere Übergänge haben Wahrscheinlichkeiten zwischen 0 und 1, die empirische
Häufigkeiten widerspiegeln, die je nach Kontext variieren können:

\begin{itemize}
    \item \textbf{KBG → VBG (0,667)} vs. \textbf{KBG → VBBd (0,333)}: In zwei
    Dritteln der Fälle wird ein Kunden-Gruß erwidert; in einem Drittel antwortet
    der Verkäufer direkt mit einer Nachfrage (Überspringen des reziproken Grußes).
    
    \item \textbf{KBBd → VBBd (0,667)} vs. \textbf{KBBd → VAA (0,167)} vs.
    \textbf{KBBd → VBA (0,167)}: Die meisten Kunden-Bedarfe werden von
    Verkäufer-Nachfragen gefolgt, manchmal aber auch direkt von Abschluss
    (sofortiger Kauf) oder Reaktion (beratende Antwort).
    
    \item \textbf{VBBd → KBA (0,444)} vs. \textbf{VBBd → VAA (0,222)} vs.
    \textbf{VBBd → KBBd (0,222)} vs. \textbf{VBBd → KAA (0,111)}: Verkäufer-
    Nachfragen haben die vielfältigsten Folgen – Antworten, Abschlüsse,
    Bedarfsschleifen (Upselling) oder Kunden-Abschlüsse (vorzeitiger Ausstieg).
    
    \item \textbf{KBA → VBA (0,5)} vs. \textbf{KBA → VAA (0,5)}: Kunden-Antworten
    führen mit gleicher Wahrscheinlichkeit zu Verkäufer-Reaktionen oder direkten
    Abschlüssen. Dies deutet auf einen strategischen Wahlpunkt hin.
    
    \item \textbf{VAA → KAA (0,857)} vs. \textbf{VAA → KAV (0,143)}: Die meisten
    Verkäufer-Abschlüsse werden von Kunden-Abschlüssen gefolgt; selten direkt
    von Kunden-Verabschiedung (abgekürzter Abschluss).
    
    \item \textbf{KAA → VAV (0,75)} vs. \textbf{KAA → VBG (0,25)}: Drei Viertel
    der Kunden-Abschlüsse führen zu Verkäufer-Verabschiedungen; ein Viertel
    führt zu einem Neustart des Gesprächs (neue Begrüßung). Die Neustart-Option
    tritt auf, wenn ein neuer Kunde eintrifft oder derselbe Kunde einen
    zusätzlichen Kauf tätigt.
\end{itemize}

\subsection{Die Upselling-Schleife}

Der Übergang \textbf{VBA → KBBd (0,5)} verdient besondere Aufmerksamkeit.
Diese Schleife – von Verkäufer-Reaktion zurück zu Kunden-Bedarf – ist die
grammatische Repräsentation des \textit{Upsellings}. Wenn ein Verkäufer
"Sonst noch etwas?" (VBA) sagt, antwortet der Kunde oft mit einem zusätzlichen
Bedarf (KBBd).

Entscheidend ist, dass diese Schleife nur in der Hälfte der Fälle auftritt
(0,5). Die andere Hälfte führt zu Abschluss (VAA, 0,25) oder Erkundigung
(KAE, 0,25). Dies deutet darauf hin, dass Upselling weder obligatorisch noch
selten ist – es ist eine strategische Option, die Verkäufer einsetzen können,
und Kunden können sie annehmen oder abweisen.

\subsection{Die Neustart-Option}

Der Übergang \textbf{KAA → VBG (0,25)} ist besonders interessant.
Ein Kunden-Abschluss (KAA) wird normalerweise von einer Verabschiedung (VAV,
0,75) gefolgt, aber in einem Viertel der Fälle folgt eine Verkäufer-Begrüßung
(VBG). Dies zeigt, dass Gespräche nach dem Abschluss neu gestartet werden
können – zum Beispiel, wenn ein neuer Kunde eintrifft (wie in T2 und T5)
oder wenn derselbe Kunde einen zusätzlichen Kauf tätigt.

\subsection{Die Ankunft eines neuen Kunden}

Der Übergang \textbf{KAV → KBBd (0,5)} zeigt, dass eine Kunden-Verabschiedung
"zurückgenommen" werden kann, wenn ein neuer Kunde eintrifft. In T5 folgt
auf eine Verabschiedung (KAV) sofort ein neuer Kunden-Bedarf (KBBd). Dies
ist der einzige Übergang, der die ansonsten strikte Phasenabfolge durchbricht.

\section{Konstitutive Regeln vs. statistische Regularitäten}

\subsection{Die Unterscheidung}

Das ARS-Framework hält eine strikte Unterscheidung zwischen zwei
Beschreibungsebenen ein:

\begin{enumerate}
    \item \textbf{Konstitutive Regeln}: Strukturelle Beschränkungen, die
    definieren, was als wohlgeformte Sequenz gilt. Diese sind binär (eine
    Sequenz entspricht entweder den Regeln oder nicht). Im Korpus sind
    Übergänge mit Wahrscheinlichkeit 1,0 Kandidaten für konstitutive Regeln.
    
    \item \textbf{Statistische Regularitäten}: Empirische Häufigkeiten von
    Übergängen. Diese sind probabilistisch und können aktualisiert werden,
    wenn neue Daten eintreffen. Sie spiegeln wider, was \textit{tatsächlich}
    vorkommt, nicht was \textit{notwendigerweise} vorkommen muss.
\end{enumerate}

\begin{table}[H]
\centering
\caption{Aus dem Korpus identifizierte konstitutive Regeln}
\label{tab:constitutive}
\begin{tabular}{@{} l l l @{}}
\toprule
\textbf{Regel} & \textbf{Wahrscheinlichkeit} & \textbf{Interpretation} \\
\midrule
VBG → KBBd & 1,0 & Verkäufer-Gruß muss von Kunden-Bedarf gefolgt sein \\
KAE → VAE & 1,0 & Kunden-Erkundigung muss beantwortet werden \\
VAE → KAA & 1,0 & Information muss zum Abschluss führen \\
VAV → KAV & 1,0 & Verabschiedungen sind reziprok \\
\bottomrule
\end{tabular}
\end{table}

\subsection{Von der Statistik zur Konstitutivität}

Eine statistische Regularität kann durch methodologische Entscheidung zu einer
konstitutiven Regel werden. Der Übergang \texttt{VBG → KBBd (1,0)} variierte
im Korpus nie. Wir könnten ihn als konstitutive Regel behandeln: "In
Verkaufsgesprächen muss auf einen Verkäufer-Gruß eine Kunden-Bedarfsäußerung
folgen." Dies verwandelt eine empirische Beobachtung in eine normative
Beschränkung.

Diese Entscheidung ist jedoch nicht automatisch. Sie erfordert methodologische
Reflexion: Ist dies wirklich ein notwendiges Merkmal des Interaktionsformats,
oder könnte es in anderen Kontexten variieren? Der ARS-Ansatz behandelt
Übergänge zunächst als statistisch, bis Gegenbeispiele eine Revision der
strukturellen Grammatik erzwingen.

\subsection{Die methodologische Bedeutung}

Diese Flexibilität ist entscheidend für XAI und neuro-symbolische Integration.
Sie erlaubt dem Analysten:

\begin{enumerate}
    \item Mit rein statistischem Lernen zu beginnen (keine vorherigen
    Beschränkungen).
    \item Übergänge zu identifizieren, die in den Daten nie variieren.
    \item Sie nach methodologischer Reflexion zu konstitutiven Regeln zu erheben.
    \item Das resultierende hybride System für Vorhersage und Erklärung zu
    nutzen.
\end{enumerate}

Das System bewegt sich so von rein empirischer Mustererkennung (System 1)
zu regelbasierter Erklärung (System 2) – ein Wechsel, der Kahnemans
Unterscheidung zwischen schnellem und langsamem Denken widerspiegelt
\citep{kahneman2011thinking}.

\section{Hin zur neuro-symbolischen Implementierung}

\subsection{Die Dual-Dynamics-Architektur}

Die ARS-Grammatik legt nahe eine Dual-Dynamics-Architektur für die
neuro-symbolische Implementierung:

\begin{enumerate}
    \item \textbf{Symbolische Komponente} (System 2): Die Grammatikregeln,
    einschließlich sowohl konstitutiver Regeln als auch statistischer
    Wahrscheinlichkeiten. Diese Komponente ist inspizierbar, falsifizierbar
    und erklärbar.
    
    \item \textbf{Neuronale Komponente} (System 1): Ein neuronales Netz,
    das Übergangswahrscheinlichkeiten aus Daten lernt und nächste Symbole
    vorhersagt. Diese Komponente ist schnell, musterbasiert und skalierbar.
    
    \item \textbf{Hybride Integration}: Die neuronale Komponente liefert
    schnelle Vorhersagen; die symbolische Komponente validiert sie gegen
    konstitutive Regeln und liefert Erklärungen.
\end{enumerate}

\begin{lstlisting}[caption=Pseudocode für neuro-symbolische ARS]
# Pseudocode: ARS 5.0 Dual-Dynamics-Architektur

class ARSNeuroSymbolicSystem:
    
    def __init__(self):
        # Symbolische Komponente
        self.grammar = ARSGrammar()           # PCFG mit Wahrscheinlichkeiten
        self.counts = zero_matrix(12, 12)    # Übergangszähler
        
        # Neuronale Komponente
        self.neural_network = NeuralNetwork( input_dim=12, hidden=64, output_dim=12 )
        
        # Konstitutive Regeln (harte Beschränkungen)
        self.constitutive_rules = {
            (KBG, VBG): True, (VBG, KBBd): True,
            (KAE, VAE): True, (VAV, KAV): True, (KAV, VAV): True
        }
    
    def update_symbolic(self, from_sym, to_sym):
        """Schnelles, zählbasiertes Update (System 2)"""
        self.counts[from_sym][to_sym] += 1
        row_sum = sum(self.counts[from_sym])
        self.grammar.probs[from_sym] = self.counts[from_sym] / row_sum
    
    def update_neural(self, from_sym, to_sym):
        """Langsames, gradientenbasiertes Update (System 1)"""
        loss = cross_entropy( self.neural_network(from_sym), to_sym )
        loss.backward()
        optimizer.step()
    
    def predict_next(self, from_sym):
        """Hybride Vorhersage"""
        neural_probs = self.neural_network(from_sym)
        symbolic_probs = self.grammar.probs[from_sym]
        return 0.5 * neural_probs + 0.5 * symbolic_probs
    
    def is_valid(self, from_sym, to_sym):
        """Prüfung konstitutiver Regeln"""
        return self.constitutive_rules.get((from_sym, to_sym), True)
\end{lstlisting}

\subsection{Erklärbarkeit durch Beweisbäume}

Die symbolische Komponente ermöglicht Erklärbarkeit durch Beweisbäume.
Für jede wohlgeformte Sequenz kann die Grammativ eine Ableitung produzieren:

\begin{lstlisting}[caption=Beweisbaum für eine wohlgeformte Sequenz]
well_formed([KBG, VBG, KBBd])
    ← start(KBG)                               [1.0]
    ← transition(KBG, VBG)                     [0.667]
    ← well_formed([VBG, KBBd])
      ← transition(VBG, KBBd)                  [1.0]
      ← well_formed([KBBd])
        ← start(KBBd)                          [0.0]
Wahrscheinlichkeit: 0.667 × 1.0 × 0.0 = 0.0
\end{lstlisting}

Dieser Beweisbaum macht das Schließen transparent. Jeder Schritt wird durch
eine Regel gerechtfertigt, und jede Regel hat eine explizite Wahrscheinlichkeit.

\subsection{Implementierungsoptionen}

Die ARS-Grammatik kann in verschiedenen neuro-symbolischen Frameworks
implementiert werden:

\begin{table}[H]
\centering
\caption{Implementierungsoptionen für ARS 5.0}
\label{tab:options}
\begin{tabular}{@{} l l l @{}}
\toprule
\textbf{Framework} & \textbf{Sprache} & \textbf{Am besten für} \\
\midrule
DeepProbLog & Prolog/Python & Forschung, Erklärbarkeit \\
PyTorch & Python & Flexibilität, Prototyping \\
Flux.jl & Julia & Wissenschaftliches Rechnen, Performance \\
Candle & Rust & Produktion, Edge-Computing \\
\bottomrule
\end{tabular}
\end{table}

Jedes Framework hat seine eigenen Stärken, aber alle können dieselbe
Dual-Dynamics-Architektur implementieren.

\section{Fazit}

Dieser Beitrag hat die empirische Grammatik von acht Marktgesprächen
rekonstruiert, von den Roh-Transkripten bis zu einer empirisch optimierten
probabilistischen Grammatik. Die Analyse ergab drei Haupterkenntnisse:

\begin{enumerate}
    \item \textbf{Empirisch}: Die optimierte Grammatik erreicht eine hohe
    Korrelation mit den Daten (\(r = 0,925\)) und offenbart sowohl konstitutive
    Regeln (Übergänge mit Wahrscheinlichkeit 1,0) als auch statistische
    Regularitäten.
    
    \item \textbf{Interpretativ}: Die Upselling-Schleife (VBA → KBBd, 0,5) und
    die Neustart-Option (KAA → VBG, 0,25) heben strategische Wahlmöglichkeiten
    in Verkaufsinteraktionen hervor. Die Verabschiedungs-Zurücknahme (KAV → KBBd, 0,5)
    zeigt, wie neue Kunden in laufende Interaktionen eintreten können.
    
    \item \textbf{Methodologisch}: Die Unterscheidung zwischen konstitutiven
    Regeln und statistischen Regularitäten bietet eine Grundlage für erklärbare,
    falsifizierbare und anpassungsfähige neuro-symbolische Systeme.
\end{enumerate}

Die hier präsentierte Grammatik ist kein statisches Artefakt. Sie kann
aktualisiert werden, wenn neue Daten eintreffen (statistische Plastizität),
und revidiert werden, wenn strukturelle Anomalien erkannt werden (strukturelle
Stabilität). In diesem Sinne ist sie ein lebendiges neuro-symbolisches Modell –
eine empirische Grammatik, die lernt, während sie erklärbar bleibt.

\newpage
\begin{thebibliography}{99}

\bibitem[Kahneman(2011)]{kahneman2011thinking}
Kahneman, D. (2011). \textit{Thinking, Fast and Slow}. Farrar, Straus and Giroux.

\bibitem[Koop(1994)]{koop1994scheme}
Koop, P. (1994). \textit{Grammatikinduktion empirisch gesicherter 
Verkaufsgespräche}. Scheme-Quellcode.

\bibitem[Koop(2026)]{koop2026deep}
Koop, P. (2026). \textit{Von Scheme zu DeepProbLog: Die ARS als methodologischer
Bauplan für moderne neuro-symbolische Programmierung}. the-last-freedom.org.

\bibitem[Manhaeve et al.(2018)]{manhaeve2018deepproblog}
Manhaeve, R., Dumancic, S., Kimmig, A., Demeester, T., \& De Raedt, L. (2018).
DeepProbLog: Neural probabilistic logic programming. \textit{Advances in
Neural Information Processing Systems}, 31.

\bibitem[Oevermann et al.(1979)]{oevermann1979methodology}
Oevermann, U., Allert, T., Konau, E., \& Krambeck, J. (1979). Die Methodologie
einer objektiven Hermeneutik. In H.-G. Soeffner (Hrsg.), \textit{Interpretative
Verfahren in den Sozial- und Textwissenschaften} (S. 352-434). Metzler.

\end{thebibliography}

\end{document}